\section{Experiments}
\label{sec:exp}

\subsection{Data, targets and uncertainty}
\label{sec:data}
We use two public score sources, both frozen snapshots; checksums and the full provenance are in Appendix~\ref{app:repro}. (i) The \textbf{Open LLM Leaderboard} dataset \texttt{open-llm-leaderboard/contents} (commit \texttt{9c09a7ca}, last updated 2025-03-20): six tasks (IFEval, BBH, MATH Lvl~5, GPQA, MuSR, MMLU-PRO), scores in the leaderboard's \emph{normalised} form, whose mean is the leaderboard's own rule; after removing duplicate model names and one flagged submission, $M=4496$ models. (ii) The \textbf{HELM} MMLU release \texttt{v1.13.0} (2025-01-10): exact-match scores on the $n=57$ MMLU subjects for the $M=79$ models that have all of them. Every model is a target, and we report results by the target's rank under the plain mean (groups 2, 3--5, 6--10, 11--30, 31--100, 101--500, $>$500; HELM has no model beyond rank 79). A group's number is the median over its targets, with a 95\% percentile-bootstrap interval over targets ($B=1000$, seed 0); a target that is infeasible under a cap counts as $+\infty$ in that median and we give the infeasible share, and ``over feasible'' medians are labelled as such. Costs of $0$ mean the target is already first under that rule. Score rules are solved for all targets and all caps. Ordinal rules are expensive (mixed-integer programs) and are solved on a rank-stratified sample of 136 OLLB targets (ranks 2--21 plus 150 log-uniform draws in rank, seed 0) and all 78 non-leading HELM models; population statements about the sample use stratum weights and are flagged. Exact Copeland, Schulze and Kemeny rankings were not computed.

\textbf{Caps.} We show \texttt{ceil} (reference upper bound on what is possible), \texttt{emp95} and \texttt{gain10}. The empirical cap \texttt{emp95} is the 95th percentile of the fine-tune-minus-base score difference per OLLB task, measured on 540 base/fine-tune pairs (154 labelled chat models, 346 fine-tuned on domain-specific datasets, 40 other); it exists only for OLLB. Its values in points are 46.0 (IFEval), 7.3 (BBH), 21.3 (MATH Lvl~5), 3.8 (GPQA), 7.2 (MuSR) and 3.2 (MMLU-PRO), the same for every target; the large IFEval and MATH values may reflect general instruction tuning rather than training on related data. Without the 154 chat-typed pairs they fall to 44.6 and 19.4 and without chat-named pairs as well they are 55.7 and 18.8 (Appendix~\ref{app:checks}); the self-declared type does not identify instruction conversion, so \texttt{emp95} is a cap on the large gains of fine-tunes and not a measurement of targeted training. \texttt{front2} (overtake the best rival on a task by at most 2 points) and the remaining caps are in the appendix.

\subsection{Results}
\label{sec:results}
\paragraph{Closed forms and solvers agree.}
For the mean, minimum (both datasets) and median (HELM; $n$ odd) the closed forms of Propositions~\ref{prop:mean}--\ref{prop:median}, coded independently of the solvers, agree with the solver output on every instance (largest absolute difference in $\LL$ $2.3\times10^{-13}$; feasibility and $K$ agree in 100\% of 35{,}960 OLLB and 468 HELM mean instances, and in the other checked cells), and on the 95{,}181 OLLB and 2{,}967 HELM instances of all rules and caps $\Linf\le \LL\le n\Linf$ and $1\le K\le n$ hold. For these three rules the solver is itself a closed-form or greedy routine, so the comparison shows that two implementations agree; the solvers, including the mixed-integer ones, were in addition checked against brute-force enumeration on small random instances. The minimum rule does \emph{not} need all tasks (Proposition~\ref{prop:min}): it needs the tasks below the best rival's minimum ($K=2$ in the median at ranks 6--500 on OLLB; 2, 7 and 17 at ranks 6--10, 11--30, 31--100 on HELM).

\begin{table}[t]
\centering\small
\caption{Median minimum $\LL$ (total points) [95\% bootstrap interval] over the \emph{feasible} targets of each rank group, per rule and cap; ``$a/b$ feas.'' means that $a$ of the group's $b$ targets can become first under the cap (cells without it: all can), ``all inf.'' that none can; $^\dagger$: median over fewer than five feasible targets, its interval is unreliable. The targets are all models of the leaderboard (a census), so the interval shows how much the median depends on which models are in the group, not sampling from a larger population; medians of different caps are over different feasible subsets. The group $>$500 is 89\% of the targets; under the binding caps it is almost wholly infeasible. Ordinal rules ($^*$) use the stratified sample (smaller $b$). Open LLM Leaderboard; the corresponding table for $K$ and for HELM MMLU are in Appendix~\ref{app:tables}.}
\label{tab:l1ollb}
\resizebox{\linewidth}{!}{\input{../figtab/tables/tab_L1_ollb.tex}}
\end{table}

\paragraph{The rule changes the cost, and the ordering of rules depends on the cost measure and the dataset.}
Figure~\ref{fig:rank} and the first numbers below are for the uncapped case. One might expect the rules that need breadth (median, minimum) to need the most investment; in points the opposite holds, in tasks it does not. In total points ($\LL$), averaging rules are the \emph{most} expensive: at ranks 101--500 on OLLB the median $\LL$ is 75 points for the mean, 55 for the trimmed mean and 51 for the geometric mean, against 29 for the median rule and 12 for the minimum; at ranks 31--100 on HELM it is 1021 for the mean and 992 for the geometric mean against 207 for the median rule, 200 for the minimum and 72 for plurality. The reason is Proposition~\ref{prop:mean}: the mean must make up $n$ times the gap in mean score, while median and minimum only have to lift specific entries to the best rival's median or minimum. In tasks touched ($K$), the order reverses for the mean: its median $K$ is 1 for every OLLB group up to ranks 101--500 (one sub-dataset with enough headroom suffices), while on HELM at ranks 31--100 the median rule needs 26 subjects (21 for the mean, 17 for the minimum) and the ordinal rules need the most: Borda 51 of 57, the Condorcet requirement 29, plurality 9. No rule is cheapest on every measure. The order of rules by median $\LL$ is also not fixed. Across caps on the same dataset the agreement is only moderate on both datasets (Figure~\ref{fig:robust}a; Kendall $\tau$ between the rule orders of two caps, for rank groups of at least five targets, 6--10 to 31--100: median 0.45, minimum 0.22 on HELM, 40 pairs; median 0.50, minimum $-0.52$ on OLLB, 77 pairs, and 0.15 at ranks 101--500, 21 pairs), it differs further between the two datasets (median 0.36, range $-0.50$ to 0.86, 131 pairs), and under evaluation-noise and task-bootstrap resampling the mean rank of a rule moves by one to several places (Figure~\ref{fig:robustfull}b,c; mean $\tau$ to the median ordering 0.74--0.77 on HELM \texttt{gain10}, 0.75 under noise and 0.35 under task bootstrap on OLLB). We therefore report the measured orders with their intervals, not a universal order; what is invariant are the closed forms.

\begin{figure}[!t]
\centering
\includegraphics[width=\linewidth]{../figtab/figures/fig1_investment_vs_rank.pdf}
\caption{Minimum investment to reach first place against the target's rank group under the plain mean, cap \texttt{ceil} (every task may be raised to the scale maximum, so all targets are feasible). Left: total points $\LL$ (log scale); right: tasks touched $K$. Top: Open LLM Leaderboard (4496 models, 6 tasks); bottom: HELM MMLU v1.13.0 (79 models, 57 subjects). Points are group medians, bars 95\% bootstrap intervals over targets; dotted lines mark the ordinal rules, which on OLLB use the stratified sample of 136 targets. Ranks 2 and 3--5 (one and three targets) are in the tables of the appendix.}
\label{fig:rank}
\end{figure}

\paragraph{Rank, and infeasibility under realistic caps.}
At \texttt{ceil}, the median $\LL$ of the largest-rank group is the largest for every rule (Figure~\ref{fig:rank}), and the group medians grow with the rank group for the arithmetic, geometric and trimmed mean on OLLB and for all rules on HELM except plurality, whose first two groups hold one and three targets. They do not for the median rule, the minimum rule, the harmonic mean and the power mean $p=-5$ on OLLB: for the minimum, the group medians from rank 2 to $>$500 are 3.7, 3.1, 13.1, 9.7, 13.8, 12.3 and 46.5 points. Under the mean, $\LL=n\cdot(\text{gap}+\varepsilon)$ (Proposition~\ref{prop:mean}), so its growth with rank is that of the gap to the leader in the leaderboard's score distribution, not an effect of the rule. As a descriptive summary only, the log--log regression of $\LL$ on rank (feasible targets, \texttt{ceil}) has slopes from 0.44 (mean) to 0.65 ($p=-5$) on OLLB and 1.07 (geometric mean) to 2.17 (median) on HELM; a power law has the smallest relative error in 25 of 32 OLLB and 16 of 18 HELM group-level fits, but a linear form has the higher $R^2$ on the $\LL$ scale in 14 of 32 and 17 of 18 (HELM: four or five groups; Appendix~\ref{app:checks}), so ``roughly a sub-linear power of the rank'' on OLLB is a description, not a demonstrated law; the slopes of the two leaderboards differ unexplained. With \texttt{emp95} the mean rule is feasible for all targets up to rank 100, 59\% at ranks 101--500 and none beyond rank 500 on OLLB, while the minimum rule is feasible for 20\% of ranks 6--10 and 6\% of ranks 101--500; the share of all OLLB targets that can become first under \texttt{emp95} is 7.5\% for the mean and 1.4\% for the minimum (\texttt{gain10}: 2.5\% and 9.9\%; Figure~\ref{fig:caps}). Which rule is easiest to topple therefore flips with the cap.

\begin{figure}[!t]
\centering
\includegraphics[width=\linewidth]{../figtab/figures/main_feasible_vs_cap.pdf}
\caption{Share of all targets that can reach first place, by cap and rule (bars: percentile bootstrap 95\% intervals over targets, within strata for the sampled ordinal rules; zero shares are drawn at the axis floor; caps are categorical, roughly from loosest to tightest). (a) OLLB (4495 targets, 89\% of them ranked below 500): 1.4--7.5\% at \texttt{emp95}, 1.2--9.9\% at \texttt{gain10}, 14--64\% at \texttt{gain20}, depending on the rule. (b) HELM MMLU (78 targets): 18--73\% at \texttt{gain10}, 6--37\% at \texttt{gain5}. Every target is feasible at \texttt{ceil} and \texttt{front2}. Ordinal rules (dotted) on OLLB: 136-target stratified sample re-weighted by stratum size. By rank group: Figure~\ref{fig:capsrank}.}
\label{fig:caps}
\end{figure}

\paragraph{Targeted versus uniform investment.}
Comparing the targeted cost $\LL$ with the cost $n\Linf$ of a uniform improvement by the same largest jump gives ratios of 1.00 (OLLB) and 0.93--0.99 (HELM) for the mean. This follows from Proposition~\ref{prop:mean} (the mean needs $n$ times the gain in mean score wherever the points go) and is not a finding; targeting saves tasks, not points. Under binding caps the mean's ratio falls to group medians of 0.42--0.89 on OLLB with \texttt{emp95} and 0.61--0.70 on HELM with \texttt{front0} (0.42--1.00 over all caps), but since $n\Linf$ ignores the cap, this mixes the effect of the cap on $\Linf$ with that of targeting; we did not compute a cap-respecting uniform raise. For the other rules the ratios are far below one (ranges over rank groups at \texttt{ceil}, OLLB and HELM): 0.33--0.51 and 0.09--0.24 for the median, 0.26--0.43 and 0.03--0.13 for the minimum, 0.26--0.29 and 0.09--0.13 for plurality.

\begin{figure}[!t]
\centering
\includegraphics[width=\linewidth]{../figtab/figures/main_targeted_vs_uniform.pdf}
\caption{Targeted versus uniform investment: ratio $\LL/(n\Linf)$ of the cheapest targeted investment to the cost of improving all $n$ tasks by the same largest jump, rank group 31--100 (OLLB $n=70$, HELM $n=49$; median over feasible targets, 95\% bootstrap interval), at \texttt{ceil} and at one binding cap per dataset (OLLB \texttt{emp95}, HELM \texttt{front0}). A ratio of one (dashed) means that targeting saves no points; $\LL\le n\Linf$ always holds, and $n\Linf$ ignores the cap. Ordinal rules on OLLB: stratified sample; under \texttt{emp95} the ratio is over the feasible targets only (Figure~\ref{fig:caps}).}
\label{fig:uniform}
\end{figure}

\paragraph{Conditions: cap, number of tasks, data slices, ties.}
The cap is the assumption on which the answer depends most (Figures~\ref{fig:caps} and~\ref{fig:capsrank}): a runner-up needs a median of 4.8 points under the mean at \texttt{ceil} on OLLB, yet beyond rank 500 no target can become first under the mean with \texttt{emp95} or \texttt{gain10}, and at ranks 101--500 the shares are 59\% and 4\% (Figure~\ref{fig:capsrank}). Without chat-typed pairs the shares under \texttt{emp95} are 6.7\% (mean), 2.7\% (median) and 1.7\% (minimum), against 7.5, 3.1 and 1.4\%; a cap on the share of the headroom a fine-tune closed gives 1.4, 1.7 and 0.8\%, and a flat quarter of the headroom 12.2, 9.4 and 100\% (Table~\ref{tab:capvar}). The number of tasks matters: on random subsets of $n$ HELM subjects (Figure~\ref{fig:robust}b) the log--log slope of the median $\LL$ in $n$ is 0.55--1.06 for the mean but $-0.14$ to 0.06 for median and minimum (five subset sizes, one benchmark family). Removing merged models or keeping the top 500 leaves the rule order close to the main one (Kendall $\tau$ 0.81--1.00); raw instead of normalised OLLB scores changes it more ($\tau$ 0.71 at \texttt{ceil}, 0.52 at \texttt{gain10}; Figure~\ref{fig:robustfull}d). Raising the tie margin from $10^{-3}$ to 0.1 points moves the median $\LL$ by 0.6--3.3\% in the 20 checked rule--cap cells, and counting ties as losses leaves it unchanged, but individual targets can change by factors up to 6.0 for Borda and 2.4 for the mean (Table~\ref{tab:eps}).

\begin{figure}[!t]
\centering
\includegraphics[width=\linewidth]{../figtab/figures/main_robustness.pdf}
\caption{Instability of the rule order and robustness. (a) Kendall $\tau$ between the rule orders (by median $\LL$ in the rank group; at least seven common rules) of every two caps of the same dataset: one dot per pair of caps (OLLB 28 pairs, 21 for ranks 31--100; HELM 15, 10), bar: median. (b) HELM, median $\LL$ at \texttt{ceil} against the number $n$ of subjects in a random subset (50 draws per $n$), target ranked fifth in the subset. (c) Mean Kendall $\tau$ of the per-draw rule order with the order of medians under evaluation noise and under task bootstrap (100 draws; \texttt{gain10}; the ten best models after the leader as targets). Intervals and per-rule ranks: Figure~\ref{fig:robustfull}.}
\label{fig:robust}
\end{figure}
